\begin{codeblock}
namespace std {
  template<class T>
  class @\libglobal{range}@ {
  public:
    // \ref{demo.range.iter}, class range::iterator
    class iterator {
    public:
      void advance();

      int get() const;

      void reset(double) = delete;

      class cursor {
      public:
        int at() const;
      };

    private:
      int steps_ = 0;
    };

    class @\exposid{sentinel}@ {
    public:
      bool reached() const;
    }; // \expos

    // \ref{demo.range.obs}, observers
    iterator begin() const;
  };
}
\end{codeblock}

\rSec3[demo.range.iter]{Class range::iterator}

\indexlibrarymember{advance}{iterator}%
\begin{itemdecl}
void advance();
\end{itemdecl}

\begin{itemdescr}
\pnum
\constraints
\tcode{steppable<T>} is \tcode{true}.

\pnum
\effects
Advances the iterator by one step.
\end{itemdescr}

\rSec3[demo.range.obs]{Observers}

\indexlibrarymember{begin}{range}%
\begin{itemdecl}
iterator begin() const;
\end{itemdecl}

\begin{itemdescr}
\pnum
\returns
An iterator to the first element.
\end{itemdescr}
